\documentclass[10pt,a4paper]{amsart}
\usepackage[margin=19mm]{geometry}
\usepackage{amsmath,amssymb,mathtools}
\usepackage[scr=rsfs,cal=euler]{mathalfa}
\usepackage{xfrac}
\usepackage[unicode,hidelinks,bookmarks=false]{hyperref}
\newcommand{\ie}{i.e.\ }
\newcommand{\cf}{cf.\ }
\newcommand{\resp}{respectively\ }
\newcommand{\Subsection}{\subsection}
\providecommand{\nobreakdash}{\nobreak\hbox{-}}
\setlength{\emergencystretch}{2em}
\multlinegap0pt
\usepackage{bm}
\usepackage{bbm}
\usepackage{dsfont}
\usepackage{xspace}
\usepackage{cancel}
\usepackage{mathtools}
\usepackage{amsthm}
\makeatletter
\@ifundefined{lem}{\newtheorem{lem}{Lem}}{}
\@ifundefined{prop}{\newtheorem{prop}[lem]{Proposition}}{}
\makeatother
\makeatletter
\newcommand{\Fcal}{\mathcal{F}}
\newcommand{\Int}{\hbox{Int}}
\newcommand{\Mcal}{\mathcal{M}}
\newcommand{\nuc}{\mathfrak{n}}
\expandafter\ifx\csname proof\endcsname\relax
\newenvironment{proof}{\noindent\textit{Proof. }}{\hfill $\Box$\linebreak}
\fi
\makeatother
\title[On compact sets possessing $q$-convex functions]{On compact sets possessing $q$-convex functions}
\author{Thomas Pawlaschyk, Nikolay Shcherbina}
\date{Source: arXiv:2505.24588v1 (CC BY 4.0)}
\begin{document}
\maketitle

\noindent\emph{Source and licence.} On compact sets possessing $q$-convex functions. Version arXiv:2505.24588v1 (CC BY 4.0), distributed under the Creative Commons Attribution 4.0 International licence, as recorded in the arXiv licence metadata for this exact version and shown on the abstract page https://arxiv.org/abs/2505.24588v1. Full rights notice: licenses/arxiv-2505.24588-CC-BY-4.0.txt.\par\smallskip

\noindent\textit{Editorial scope.} Counted unit: the Prop whose source label is \texttt{prop-qnuc-qpsc} in the article (source0.tex lines 449--450, source label \texttt{prop-qnuc-qpsc}), delivered together with the article's own complete original proof of it (source0.tex lines 452--478, one \texttt{proof} environment of 27 lines). No mathematics is written, repaired or re-derived: statement and proof are the article's own text line for line. Required context retained uncounted: prop-qpsc-caps (lines 173--186). Every definition, display and label of those lines is retained unchanged. Omitted from this package and not redistributed: 1--172, 187--448, 451--451, 479--791. Nothing inside a retained excerpt is omitted or altered: every retained body is delivered exactly as the article's lines are written, so no omission receipt of any kind accompanies this package. Citations: the retained excerpts carry no citation command at all, so no bibliography entry of the article is transcribed and no reference is redistributed. Reference adaptation: none was needed and none was made; every reference inside the delivered unit resolves inside it, so no unresolved reference or citation is produced. Printed slips kept literal: no printed word, symbol, hypothesis or number of the counted statement or of its proof was altered. Layout and class: the article's own class is \texttt{article} with packages including amsmath, amssymb, ntheorem, hyperref, verbatim, authblk, xcolor; the wrapper is the lane's pinned \texttt{amsart} editorial wrapper at 10pt on A4, which restates the theorem-like environments the retained text needs and adds the article's own macro definitions that the retained text needs (\texttt{\textbackslash Fcal}, \texttt{\textbackslash Int}, \texttt{\textbackslash Mcal}, \texttt{\textbackslash nuc}), with extra wrapper packages (bm, bbm, dsfont, xspace, cancel, mathtools). The delivered numbering is the unit's own. Assets: no retained span contains a figure environment, a graphics-inclusion command, a PDF-page inclusion, a table or a TikZ picture; no image file is copied into this package, the package declares no asset dependency and no image file is redistributed with it. Licence: version arXiv:2505.24588v1 of this article is distributed under the Creative Commons Attribution 4.0 International licence, as recorded in the arXiv licence metadata for this exact version and shown on the abstract page https://arxiv.org/abs/2505.24588v1. A full rights notice is delivered at licenses/arxiv-2505.24588-CC-BY-4.0.txt. Characters: every retained excerpt is pure ASCII, so the delivered engine, XeLaTeX with its default Latin Modern text font, reports no missing character.
\begin{prop} \label{prop-qpsc-caps} Let $\Omega$ be an open set in $\Mcal$. Then the following properties are equivalent.

\begin{enumerate}

\item \label{prop-qpsc-caps-1}  $\Omega$ is $q$-pseudoconvex.

\item \label{prop-qpsc-caps-4} For every spherical hat pair $(S,\widehat{S})$ of order $n-q+1$ such that $\Omega$ contains the spherical hat $S=S^{n-q+1}$, the set $\Omega$ also contains the filled hat $\widehat{S}=\widehat{S}^{n-q+1}$.


%For every real number $0 < r < 1$, every open neighborhood $U$ of $\widehat{\SS}_r^{n-q+1}$ and every injective holomorphic map $\Phi:U\times\Delta^{q-1}\to\Mcal$ such that $\Phi(\SS_r^{n-q+1}\times\Delta^{q-1}) \subset \Omega$ we have $\Phi(\widehat{\SS}_r^{n-q+1}\times\Delta^{q-1}) \subset \Omega$.

\end{enumerate}

\end{prop}

\begin{prop}\label{prop-qnuc-qpsc} The $q$-nucleus of $K \subsetneq \Mcal$ is $q$-pseudoconcave in $\Mcal$. Moreover, it is the maximal $q$-pseudoconcave subset of $K$. In particular, if $K$ is $q$-pseudoconcave itself, then $\nuc_q(K)=K$. 
\end{prop}

\begin{proof} If the $q$-nucleus is empty, there is nothing to show, so we assume it is non-empty. Assume that $\nuc_q(K)$ is not $q$-pseudoconcave. 
%According to Proposition~\ref{prop-qpsc-caps} there exists an open neighborhood $U$ of a filled spherical hat $\widehat{\SS}_r^{n-q+1}$ and an injective holomorphic map $\Phi:U\times\Delta^{q-1}\to\Mcal\setminus A$ such that 
According to Proposition~\ref{prop-qpsc-caps} there exists a spherical hat pair $(S,\widehat{S})$ of order $n{-}q{+}1$ such that

\begin{itemize}

\item $S \cap \nuc_q(K)=\emptyset$,

\item but $\widehat{S} \cap \nuc_q(K) \neq \emptyset$.

\end{itemize}

Pick a point $p \in \widehat{S} \cap \nuc_q(K)$. Then, in view of Lemma~3.3 and the definition of the $q$-nucleus, there exists a set $K'$ in $\Fcal_K^q$ such that

\begin{itemize}

\item $S \cap K'=\emptyset$,

\item but $\widehat{S} \cap K' \neq \emptyset$. More precisely, this intersection contains $p$.

\end{itemize}

We set $K'':= K' \setminus \Int(\widehat{S})$. It is obvious that $p \notin K''$ and that $K''$ is obtained from $K'$ by a spherical cut of order $n{-}q{+}1$. Since $K'$ is obtained from $K$ by a sequence of spherical cuts of order $n{-}q{+}1$, so is $K''$ as well. This means that $K''$ lies in $\Fcal_K^q$. Hence, by the definition of the $q$-nucleus and the choice of $p$, we have that $p \in \nuc_q(K) \subset K''$, so $p \in K''$, a contradiction. Therefore, $\nuc_q(K)$ is $q$-pseudoconcave.

Now assume that there exists a $q$-pseudoconcave set $A \subset K$. Let $K''\subsetneq\Mcal$ be a compact sets in $\Mcal$ such that $K''$ is obtained from $K$ by a spherical cut of order $n{-}q{+}1$, i.e., there exists a spherical hat pair $(S,\widehat{S})$ of order $n{-}q{+}1$ such that $S \subset \Mcal\setminus K$ and $K''=K \setminus \Int(\widehat{S})$. Since $S \subset \Mcal\setminus K \subset \Mcal\setminus A$ and $A$ is $q$-pseudoconcave, $\widehat{S} \cap A = \emptyset$ due to Proposition~\ref{prop-qpsc-caps}. Hence, $A$ lies in $K''$. Similarly, by an inductive argument, we can derive that if $K''$ is obtained from $K$ by a \emph{sequence} of spherical cuts of order $n{-}q{+}1$, then $A \subset K''$ as well. But this means that $A \subset \nuc_q(K)$ according to the definition of the $q$-nucleus of $K$. In this sense, the $q$-nucleus is maximal. It contains all $q$-pseudoconcave subsets of $K$. So if $K$ is $q$-pseudoconcave itself, then $K \subset \nuc_q(K) \subset K$. Thus, $\nuc_q(K)=K$.

\end{proof}

\end{document}
